\scanpage{72}
\mspage{10}
% archive foliation: 74
\uline{NB 1} Cette fonction est unique, \unsure{dû au fait}
que le groupe des automorphismes \deleted{conformes}
holomorphes de $\mathbb D$ agit de façon simple-
ment transitive sur l'ens des
triples \deleted{\unsure{ordonnés}} de pts distincts $(r,s,t)$ \inserted{sur $S^1=\delta\mathbb D$} qui se
suivent pour l'orientation canonique de
\deleted{$\unsure{\delta\mathbb D}$} $S^1$. \deleted{\unsure{Donc $f$}}

\uline{NB 2} La condition 3) implique qu'on a un
\unsure{homéo} canonique
\[
\varphi:\beta=\widehat{(1,\zeta)}\overset\sim\longrightarrow
\overline\beta=\widehat{(1,\overline\zeta)}
\]
tel que $\varphi(1)=1$, $\varphi(\zeta)=\overline\zeta$).

B) Cette application $\varphi$ est rectifiable. Inverse-
ment, tout \deleted{application}\inserted{homéom} rectifiable $\varphi:\beta\overset\sim\longrightarrow\overline\beta$
provient d'une situation comme en A.
De façon plus précise, \deleted{on peut}\deleted{il existe} (quitte
à \unsure{remplacer} $\overline V$ par $\mathbb D$ par le th. d'uniformi-
sation A) simplifié sans segment intérieur $\overline I$,
et normaliser la situation en \deleted{\unsure{faisant}} envoyant
$a$ en $0$, $b$ en $1$) \deleted{il existe un} Il existe
une application continue
\[
f=f_\varphi:\mathbb D\longrightarrow\mathbb D
\]
telle que

1) $f|\mathbb D^\circ$ soit \deleted{holo} injection holomorphe
